\documentclass[12pt,a4paper]{article}
\usepackage[utf8]{inputenc}
\usepackage[russian]{babel}
\usepackage{hyperref}
\usepackage{xcolor}
\title{Отчет по публикациям с портала arXiv.org}
\author{Сгенерировано PolynomialDivisionApp (ICSR)}
\date{\today}
\begin{document}
\maketitle
\hrule\vspace{0.5cm}
\vspace{0.2cm}\hrule\vspace{0.2cm}
\section*{СВЕЖИЕ НАУЧНЫЕ ПУБЛИКАЦИИ И СТАТЬИ С МЕЖДУНАРОДНОГО ПОРТАЛА arXiv}
\vspace{0.2cm}\hrule\vspace{0.2cm}
\subsection*{📄 СТАТЬЯ №1 [2014-01-17]}
\begin{description}
    \item[Название:] Integral closure of rings of integer-valued polynomials on algebras
    \item[Категория:] math.AC
    \item[Авторы:] Giulio Peruginelli, Nicholas J. Werner
    \item[Ссылка на PDF:] \href{http://arxiv.org/abs/1401.4438v1}{\color{blue}http://arxiv.org/abs/1401.4438v1}
    \item[LaTeX Ключ цитирования:] \textbf{\color{teal}\cite{Peruginelli2014\_14014438}}
    \item[Аннотация (Abstract) англ:] Let \$D\$ be an integrally closed domain with quotient field \$K\$. Let \$A\$ be a torsion-free \$D\$-algebra that is finitely generated as a \$D\$-module. For every \$a\$ in \$A\$ we consider its minimal polynomial \$μ\_a(X)\in D[X]\$, i.e. the monic polynomial of least degree such that \$μ\_a(a)=0\$. The ring \${\rm Int}\_K(A)\$ consists of polynomials in \$K[X]\$ that send elemen...
\end{description}
\vspace{0.2cm}\hrule\vspace{0.2cm}
\vspace{0.2cm}\hrule\vspace{0.2cm}
\subsection*{📄 СТАТЬЯ №2 [2025-02-26]}
\begin{description}
    \item[Название:] Some permutation polynomials via linear translators
    \item[Категория:] math.NT
    \item[Авторы:] Xuan Pang, Pingzhi Yuan, Hongjian Li
    \item[Ссылка на PDF:] \href{http://arxiv.org/abs/2502.18759v1}{\color{blue}http://arxiv.org/abs/2502.18759v1}
    \item[LaTeX Ключ цитирования:] \textbf{\color{teal}\cite{Pang2025\_250218759}}
    \item[Аннотация (Abstract) англ:] Permutation polynomials with explicit constructions over finite fields have long been a topic of great interest in number theory. In recent years, by applying linear translators of functions from \$\mathbb{F}\_{q^n}\$ to \$\mathbb{F}\_q\$, many scholars constructed some classes of permutation polynomials. Motivated by previous works, we first naturally extend the ...
\end{description}
\vspace{0.2cm}\hrule\vspace{0.2cm}
\vspace{0.2cm}\hrule\vspace{0.2cm}
\subsection*{📄 СТАТЬЯ №3 [2014-10-17]}
\begin{description}
    \item[Название:] A polynomial distribution applied to income and wealth distribution
    \item[Категория:] q-fin.GN
    \item[Авторы:] Elvis Oltean, Fedor Kusmartsev
    \item[Ссылка на PDF:] \href{http://arxiv.org/abs/1410.4866v2}{\color{blue}http://arxiv.org/abs/1410.4866v2}
    \item[LaTeX Ключ цитирования:] \textbf{\color{teal}\cite{Oltean2014\_14104866}}
    \item[Аннотация (Abstract) англ:] Income and wealth distribution affect stability of a society to a large extent and high inequality affects it negatively. Moreover, in the case of developed countries, recently has been proven that inequality is closely related to all negative phenomena affecting society. So far, Econophysics papers tried to analyse income and wealth distribution by employin...
\end{description}
\vspace{0.2cm}\hrule\vspace{0.2cm}
\vspace{0.2cm}\hrule\vspace{0.2cm}
\subsection*{📄 СТАТЬЯ №4 [2021-06-21]}
\begin{description}
    \item[Название:] Explicit Multi-point Taylor Polynomial
    \item[Категория:] math.CA
    \item[Авторы:] Andrés Gómez Arias
    \item[Ссылка на PDF:] \href{http://arxiv.org/abs/2106.11440v1}{\color{blue}http://arxiv.org/abs/2106.11440v1}
    \item[LaTeX Ключ цитирования:] \textbf{\color{teal}\cite{Arias2021\_210611440}}
    \item[Аннотация (Abstract) англ:] The multi-point Taylor polynomial, which is the general, unique and of minimum degree (\$mk+m-1\$) polynomial \$P\_{k,m}(x)\$ which interpolates a function's derivatives in multiple points is presented in its explicit form. A proof that this expression satisfies the multi-point Taylor polynomial's defining property is given. Namely, it is proven that for a k-diff...
\end{description}
\vspace{0.2cm}\hrule\vspace{0.2cm}
\vspace{0.2cm}\hrule\vspace{0.2cm}
\subsection*{📄 СТАТЬЯ №5 [2024-10-01]}
\begin{description}
    \item[Название:] The HOMFLY Polynomial of a Forest Quiver
    \item[Категория:] math.CO
    \item[Авторы:] Amanda Schwartz
    \item[Ссылка на PDF:] \href{http://arxiv.org/abs/2410.00399v2}{\color{blue}http://arxiv.org/abs/2410.00399v2}
    \item[LaTeX Ключ цитирования:] \textbf{\color{teal}\cite{Schwartz2024\_241000399}}
    \item[Аннотация (Abstract) англ:] We define the HOMFLY polynomial of a forest quiver \$Q\$ using a recursive definition on the underlying graph of the quiver. We then show that this polynomial is equal to the HOMFLY polynomial of any plabic link which comes from a connected plabic graph whose quiver is \$Q\$. We also prove a closed-form expression for the HOMFLY polynomial of a forest quiver \$Q\$...
\end{description}
\vspace{0.2cm}\hrule\vspace{0.2cm}
\vspace{0.2cm}\hrule\vspace{0.2cm}
\subsection*{📄 СТАТЬЯ №6 [2017-09-22]}
\begin{description}
    \item[Название:] Quantum mechanics with orthogonal polynomials
    \item[Категория:] quant-ph
    \item[Авторы:] A. D. Alhaidari
    \item[Ссылка на PDF:] \href{http://arxiv.org/abs/1709.07652v2}{\color{blue}http://arxiv.org/abs/1709.07652v2}
    \item[LaTeX Ключ цитирования:] \textbf{\color{teal}\cite{Alhaidari2017\_170907652}}
    \item[Аннотация (Abstract) англ:] We present a formulation of quantum mechanics based on orthogonal polynomials. The wavefunction is expanded over a complete set of square integrable basis in configuration space where the expansion coefficients are orthogonal polynomials in the energy. Information about the corresponding physical systems (both structural and dynamical) are derived from the p...
\end{description}
\vspace{0.2cm}\hrule\vspace{0.2cm}
\vspace{0.2cm}\hrule\vspace{0.2cm}
\subsection*{📄 СТАТЬЯ №7 [2014-10-14]}
\begin{description}
    \item[Название:] An econophysical approach of polynomial distribution applied to income and expenditure
    \item[Категория:] q-fin.GN
    \item[Авторы:] Elvis Oltean
    \item[Ссылка на PDF:] \href{http://arxiv.org/abs/1410.3860v2}{\color{blue}http://arxiv.org/abs/1410.3860v2}
    \item[LaTeX Ключ цитирования:] \textbf{\color{teal}\cite{Oltean2014\_14103860}}
    \item[Аннотация (Abstract) англ:] Polynomial distribution can be applied to dynamical systems in certain situations. Macroeconomic systems characterized by economic variables such as income and wealth can be modelled similarly using polynomials. We extend our previous work to data regarding income from a more diversified pool of countries, which contains developed countries with high income,...
\end{description}
\vspace{0.2cm}\hrule\vspace{0.2cm}
\vspace{0.2cm}\hrule\vspace{0.2cm}
\subsection*{📄 СТАТЬЯ №8 [2017-08-27]}
\begin{description}
    \item[Название:] Topological Tutte Polynomial
    \item[Категория:] math.CO
    \item[Авторы:] Sergei Chmutov
    \item[Ссылка на PDF:] \href{http://arxiv.org/abs/1708.08132v1}{\color{blue}http://arxiv.org/abs/1708.08132v1}
    \item[LaTeX Ключ цитирования:] \textbf{\color{teal}\cite{Chmutov2017\_170808132}}
    \item[Аннотация (Abstract) англ:] This is a survey recent works on topological extensions of the Tutte polynomial....
\end{description}
\vspace{0.2cm}\hrule\vspace{0.2cm}
\vspace{0.2cm}\hrule\vspace{0.2cm}
\subsection*{📄 СТАТЬЯ №9 [2018-08-18]}
\begin{description}
    \item[Название:] On the solution existence and stability of polynomial optimization problems
    \item[Категория:] math.OC
    \item[Авторы:] Vu Trung Hieu
    \item[Ссылка на PDF:] \href{http://arxiv.org/abs/1808.06100v6}{\color{blue}http://arxiv.org/abs/1808.06100v6}
    \item[LaTeX Ключ цитирования:] \textbf{\color{teal}\cite{Hieu2018\_180806100}}
    \item[Аннотация (Abstract) англ:] This paper introduces and investigates a regularity condition in the asymptotic sense for optimization problems whose objective functions are polynomial. Under this regularity condition, the normalization argument in asymptotic analysis enables us to see the solution existence as well as the solution stability of these problems. We prove a Frank-Wolfe type t...
\end{description}
\vspace{0.2cm}\hrule\vspace{0.2cm}
\vspace{0.2cm}\hrule\vspace{0.2cm}
\subsection*{📄 СТАТЬЯ №10 [2013-07-04]}
\begin{description}
    \item[Название:] Efficiently determining Convergence in Polynomial Recurrence Sequences
    \item[Категория:] cs.DM
    \item[Авторы:] Deepak Ponvel Chermakani
    \item[Ссылка на PDF:] \href{http://arxiv.org/abs/1307.2164v1}{\color{blue}http://arxiv.org/abs/1307.2164v1}
    \item[LaTeX Ключ цитирования:] \textbf{\color{teal}\cite{Chermakani2013\_13072164}}
    \item[Аннотация (Abstract) англ:] We derive the necessary and sufficient condition, for a given Polynomial Recurrence Sequence to converge to a given target rational K. By converge, we mean that the Nth term of the sequence, is equal to K, as N tends to positive infinity. The basic idea of our approach is to construct a univariate polynomial equation in x, whose coefficients correspond to th...
\end{description}
\vspace{0.2cm}\hrule\vspace{0.2cm}
\vspace{0.2cm}\hrule\vspace{0.2cm}
\subsection*{📄 СТАТЬЯ №11 [2011-05-03]}
\begin{description}
    \item[Название:] Polynomial representation of Fermat's Last Theorem
    \item[Категория:] math.GM
    \item[Авторы:] D. De Pedis
    \item[Ссылка на PDF:] \href{http://arxiv.org/abs/1105.0669v5}{\color{blue}http://arxiv.org/abs/1105.0669v5}
    \item[LaTeX Ключ цитирования:] \textbf{\color{teal}\cite{Pedis2011\_11050669}}
    \item[Аннотация (Abstract) англ:] We propose a new approach at Fermat's Last Theorem (FLT) solution: for each FLT equation we associate a polynomial of the same degree. The study of the roots of the polynomial allows us to investigate the FLT validity. This technique, certainly within the reach of Fermat himself, allows us infer that this is the marvelous proof that Fermat claimed to have....
\end{description}
\vspace{0.2cm}\hrule\vspace{0.2cm}
\vspace{0.2cm}\hrule\vspace{0.2cm}
\subsection*{📄 СТАТЬЯ №12 [2023-10-25]}
\begin{description}
    \item[Название:] Squarefree Values Of Polynomials
    \item[Категория:] math.GM
    \item[Авторы:] N. A. Carella
    \item[Ссылка на PDF:] \href{http://arxiv.org/abs/2310.16952v1}{\color{blue}http://arxiv.org/abs/2310.16952v1}
    \item[LaTeX Ключ цитирования:] \textbf{\color{teal}\cite{Carella2023\_231016952}}
    \item[Аннотация (Abstract) англ:] This note presents new results for the squarefree value sets of quartic polynomials over the integers....
\end{description}
\vspace{0.2cm}\hrule\vspace{0.2cm}
\vspace{0.2cm}\hrule\vspace{0.2cm}
\subsection*{📄 СТАТЬЯ №13 [2019-06-06]}
\begin{description}
    \item[Название:] Tutte Polynomial Activities
    \item[Категория:] math.CO
    \item[Авторы:] Spencer Backman
    \item[Ссылка на PDF:] \href{http://arxiv.org/abs/1906.02781v2}{\color{blue}http://arxiv.org/abs/1906.02781v2}
    \item[LaTeX Ключ цитирования:] \textbf{\color{teal}\cite{Backman2019\_190602781}}
    \item[Аннотация (Abstract) англ:] Unlike Whitney's definition of the corank-nullity generating function \$T(G;x+1,y+1)\$, Tutte's definition of his now eponymous polynomial \$T(G;x,y)\$ requires a total order on the edges of which the polynomial is a posteriori independent. Tutte presented his definition in terms of internal and external activities of maximal spanning forests. Although Tutte's o...
\end{description}
\vspace{0.2cm}\hrule\vspace{0.2cm}
\vspace{0.2cm}\hrule\vspace{0.2cm}
\subsection*{📄 СТАТЬЯ №14 [2018-05-21]}
\begin{description}
    \item[Название:] On A Polytime Factorization Algorithm for Multilinear Polynomials Over F2
    \item[Категория:] cs.DM
    \item[Авторы:] Pavel Emelyanov, Denis Ponomaryov
    \item[Ссылка на PDF:] \href{http://arxiv.org/abs/1805.08186v5}{\color{blue}http://arxiv.org/abs/1805.08186v5}
    \item[LaTeX Ключ цитирования:] \textbf{\color{teal}\cite{Emelyanov2018\_180508186}}
    \item[Аннотация (Abstract) англ:] In 2010, A. Shpilka and I. Volkovich established a prominent result on the equivalence of polynomial factorization and identity testing. It follows from their result that a multilinear polynomial over the finite field of order 2 can be factored in time cubic in the size of the polynomial given as a string. Later, we have rediscovered this result and provided...
\end{description}
\vspace{0.2cm}\hrule\vspace{0.2cm}
\vspace{0.2cm}\hrule\vspace{0.2cm}
\subsection*{📄 СТАТЬЯ №15 [2015-06-13]}
\begin{description}
    \item[Название:] Polynomial and Analytic Functors and Monads, revisited
    \item[Категория:] math.CT
    \item[Авторы:] Stanisław Szawiel, Marek Zawadowski
    \item[Ссылка на PDF:] \href{http://arxiv.org/abs/1506.04317v2}{\color{blue}http://arxiv.org/abs/1506.04317v2}
    \item[LaTeX Ключ цитирования:] \textbf{\color{teal}\cite{Szawiel2015\_150604317}}
    \item[Аннотация (Abstract) англ:] We describe an abstract 2-categorical setting to study various notions of polynomial and analytic functors and monads....
\end{description}
\vspace{0.2cm}\hrule\vspace{0.2cm}
\vspace{0.2cm}\hrule\vspace{0.2cm}
\subsection*{📄 СТАТЬЯ №16 [2004-10-05]}
\begin{description}
    \item[Название:] Integer Polynomial Optimization in Fixed Dimension
    \item[Категория:] math.OC
    \item[Авторы:] Jesús A. De Loera, Raymond Hemmecke, Matthias Köppe, Robert Weismantel
    \item[Ссылка на PDF:] \href{http://arxiv.org/abs/math/0410111v2}{\color{blue}http://arxiv.org/abs/math/0410111v2}
    \item[LaTeX Ключ цитирования:] \textbf{\color{teal}\cite{Loera2004\_math/0410111}}
    \item[Аннотация (Abstract) англ:] We classify, according to their computational complexity, integer optimization problems whose constraints and objective functions are polynomials with integer coefficients and the number of variables is fixed. For the optimization of an integer polynomial over the lattice points of a convex polytope, we show an algorithm to compute lower and upper bounds for...
\end{description}
\vspace{0.2cm}\hrule\vspace{0.2cm}
\vspace{0.2cm}\hrule\vspace{0.2cm}
\subsection*{📄 СТАТЬЯ №17 [2016-09-02]}
\begin{description}
    \item[Название:] Asymptotic For Primitive Roots Producing Polynomials
    \item[Категория:] math.GM
    \item[Авторы:] N. A. Carella
    \item[Ссылка на PDF:] \href{http://arxiv.org/abs/1609.01147v3}{\color{blue}http://arxiv.org/abs/1609.01147v3}
    \item[LaTeX Ключ цитирования:] \textbf{\color{teal}\cite{Carella2016\_160901147}}
    \item[Аннотация (Abstract) англ:] Let \$x \geq 1\$ be a large number, let \$f(x) \in \mathbb{Z}[x]\$ be a prime polynomial of degree \$\text{deg}(f)=m\$, and let \$u\ne \pm 1, v^2\$ be a fixed integer. Assuming the Bateman-Horn conjecture, an asymptotic counting function for the number of primes \$p=f(n) \leq x\$ with a fixed primitve root \$u\$ is derived in this note....
\end{description}
\vspace{0.2cm}\hrule\vspace{0.2cm}
\vspace{0.2cm}\hrule\vspace{0.2cm}
\subsection*{📄 СТАТЬЯ №18 [2016-08-17]}
\begin{description}
    \item[Название:] Approximating the Chromatic Polynomial
    \item[Категория:] cs.DM
    \item[Авторы:] Yvonne Kemper, Isabel Beichl
    \item[Ссылка на PDF:] \href{http://arxiv.org/abs/1608.04883v1}{\color{blue}http://arxiv.org/abs/1608.04883v1}
    \item[LaTeX Ключ цитирования:] \textbf{\color{teal}\cite{Kemper2016\_160804883}}
    \item[Аннотация (Abstract) англ:] Chromatic polynomials are important objects in graph theory and statistical physics, but as a result of computational difficulties, their study is limited to graphs that are small, highly structured, or very sparse. We have devised and implemented two algorithms that approximate the coefficients of the chromatic polynomial \$P(G,x)\$, where \$P(G,k)\$ is the num...
\end{description}
\vspace{0.2cm}\hrule\vspace{0.2cm}
\vspace{0.2cm}\hrule\vspace{0.2cm}
\subsection*{📄 СТАТЬЯ №19 [2026-07-10]}
\begin{description}
    \item[Название:] PCPOP.jl: A Julia package for partially commutative polynomial optimization
    \item[Категория:] quant-ph
    \item[Авторы:] Moisés Bermejo Morán, Abhishek Mishra
    \item[Ссылка на PDF:] \href{http://arxiv.org/abs/2607.09339v2}{\color{blue}http://arxiv.org/abs/2607.09339v2}
    \item[LaTeX Ключ цитирования:] \textbf{\color{teal}\cite{Morán2026\_260709339}}
    \item[Аннотация (Abstract) англ:] Here we present PCPOP, a Julia package for polynomial optimization that supports non-commutative optimization, tracial polynomial optimization, trace polynomial optimization and state polynomial optimization. PCPOP fully supports exact arithmetic computations and incorporates convenient functionalities such as algebraic reductions based on Gröbner basis meth...
\end{description}
\vspace{0.2cm}\hrule\vspace{0.2cm}
\vspace{0.2cm}\hrule\vspace{0.2cm}
\subsection*{📄 СТАТЬЯ №20 [2025-11-12]}
\begin{description}
    \item[Название:] A polynomially accelerated fixed-point iteration for vector problems
    \item[Категория:] math.NA
    \item[Авторы:] Francesco Alemanno
    \item[Ссылка на PDF:] \href{http://arxiv.org/abs/2511.09012v2}{\color{blue}http://arxiv.org/abs/2511.09012v2}
    \item[LaTeX Ключ цитирования:] \textbf{\color{teal}\cite{Alemanno2025\_251109012}}
    \item[Аннотация (Abstract) англ:] Fixed-point solvers are ubiquitous in nonlinear PDEs, yet their progress collapses whenever the Jacobian at the solution carries an eigenvalue arbitrarily close to one. We ask whether such stagnation can be removed without storing long histories or solving dense least squares. Under two assumptions -- (A1) the linearised error \$e\_n\$ is dominated by a multipl...
\end{description}
\vspace{0.2cm}\hrule\vspace{0.2cm}
\vspace{0.2cm}\hrule\vspace{0.2cm}
\newpage
\bibliographystyle{unsrt}
\bibliography{references}
\end{document}